\uuid{KJTX}
\exo7id{1942}
\titre{exo7 1942}
\auteur{gineste}
\organisation{exo7}
\datecreate{2001-11-01}
\isIndication{false}
\isCorrection{false}
\chapitre{Série numérique}
\sousChapitre{Série à  termes positifs}
\cpge{mpsi-mp2i.procedes-sommatoires-discrets.a, pcsi.series-numeriques.a, ptsi.series-numeriques.a, mpsi-mp2i.procedes-sommatoires-discrets.b, pcsi.series-numeriques.b, ptsi.series-numeriques.b, bcpst1.suites-reelles, mp-mpi.series-numeriques-vectorielles.a, pt.series-numeriques.b, bcpst2.series-reelles}
\module{Analyse}
\niveau{L1}
\difficulte{}

\contenu{
\question{
Soit $u_n$ une suite décroissante à termes positifs. On suppose $(\sum u_n)$ converge. Montrer que $$ \lim_{n \rightarrow \infty} (nu_n)=0 .$$
Indication : Encadrer $\sum_{p+1}^n u_k$ pour $n>p$. Puis revenir aux définitions des limites avec les epsilons.
}
}
